\input{../tex-inputs/latex-leseansicht-vorspann}

\section[Arthur Schnitzler an Gustav Schwarzkopf, 27. 11. 1893]{L04106 Arthur Schnitzler an Gustav Schwarzkopf, 27. 11. 1893}
\nopagebreak\mylabel{L04106v}
\rehead{ }\normalsize\beginnumbering\briefempfaengerindex{Schwarzkopf, Gustav@\textsc{Schwarzkopf, Gustav}!zzzSchnitzler, Arthur@\emph{von Arthur Schnitzler}!1893-11-271@{27. 11. 1893}|(be}
\toendnotes[C]{\smallbreak\pagebreak[2]}
\correspDesc{Versand  durch Arthur Schnitzler am 27. 11. 1893 in Wien
\newline{}Erhalt  durch Gustav Schwarzkopf im Zeitraum [27. 11. 1893 – 28. 11. 1893] in Wien}\toendnotes[C]{\smallbreak}
\Standort{CUL, Schnitzler, B 96.}
\physDesc{Brief, 1 Blatt, 2 Seiten, 682 Zeichen (Blatt mir seitlichen Abrissspuren)
\newline{}Handschrift: Bleistift, deutsche Kurrent}\toendnotes[C]{\smallbreak}
\pstart
           \noindent{}{\pb}Mein lieber Freund, ich brauche Sie dringend zu einer \label{K_L04106-1v}\edtext{Probe, womöglich Mittwoch\eventindex{Volkstheater@\textbf{Volkstheater}!Probe von Das Märchen, 29.11.1893@Probe von Das Märchen, 29.11.1893|pwv}}{\lemma{\textnormal{\emph{Probe, … Mittwoch}}}\Cendnote{\textnormal{Vgl. A. S.: \emph{Kulturveranstaltungen}, 29. 11. 1893.
                  }}}\label{K_L04106-1}. Es handelt ſich
               um den unglückſeligen Moritzki\pwindex{Schnitzler, Arthur 15.\,5.\,1862 Wien – 21.\,10.\,1931 ebd.@\textsc{Schnitzler, Arthur} (15.\,5.\,1862 Wien – 21.\,10.\,1931 ebd.), \emph{Schriftsteller, Mediziner}!Märchen. Schauspiel in drei Aufzügen@\strich\emph{Das Märchen. Schauspiel in drei Aufzügen}|pwv}, über den ich tagtäglich
               mit Kadelburg\pwindex{Kadelburg, Heinrich 14.\,2.\,1856 Budapest – 13.\,7.\,1910 Marienbad@\textsc{Kadelburg, Heinrich} (14.\,2.\,1856 Budapest – 13.\,7.\,1910 Marienbad), \emph{Schriftsteller, Regisseur, Schauspieler}|pw} zu{ }ſtreiten habe. Wir haben uns heute
               geeinigt, \uline{Ihnen} die letzte Entſcheidung zu
               überlaſſen, und ich bitte Sie ſehr, mir dieſen Liebesdienſt nicht abzuſchlagen.
               Wollen Sie mir mit 2 Worten mittheilen, ob Sie Mittwoch Vormittag für
               dieſe Sache Zeit haben? Sie laſſen da{\geminationn}{\pb}einfach \textsc{Kadelburg\pwindex{Kadelburg, Heinrich 14.\,2.\,1856 Budapest – 13.\,7.\,1910 Marienbad@\textsc{Kadelburg, Heinrich} (14.\,2.\,1856 Budapest – 13.\,7.\,1910 Marienbad), \emph{Schriftsteller, Regisseur, Schauspieler}|pw}} herausrufen; – der 3. Akt\pwindex{Schnitzler, Arthur 15.\,5.\,1862 Wien – 21.\,10.\,1931 ebd.@\textsc{Schnitzler, Arthur} (15.\,5.\,1862 Wien – 21.\,10.\,1931 ebd.), \emph{Schriftsteller, Mediziner}!Märchen. Schauspiel in drei Aufzügen@\strich\emph{Das Märchen. Schauspiel in drei Aufzügen}|pw} dürfte zwiſchen ½ 
               12 u 12 begi{\geminationn}en. We{\geminationn} Sie
               der ganzen Probe\eventindex{Volkstheater@\textbf{Volkstheater}!Probe von Das Märchen, 29.11.1893@Probe von Das Märchen, 29.11.1893|pw} beiwohnen wollen, iſt es mir gewiſs ein
               großer, – Ihnen dagegen wahrſcheinlich gar kein Vergnügen?–\pend
           
\pstart
           Herzlich grüßt{\\[\baselineskip]} Ihr \spacefill\mbox{ArthSchn}\pend
           \leftskip=0em{}
\pstart
           27. 11. 93.\pend
           
\pstart
           Hrn. Gustav Schwarzkopf\pend
           \selectlanguage{ngerman}\endnumbering\briefempfaengerindex{Schwarzkopf, Gustav@\textsc{Schwarzkopf, Gustav}!zzzSchnitzler, Arthur@\emph{von Arthur Schnitzler}!1893-11-271@{27. 11. 1893}|)be}\mylabel{L04106h}
\begin{anhang}
\end{anhang}\newcommand{\dateiname}{L04106}\newcommand{\titel}{Arthur Schnitzler an Gustav Schwarzkopf, 27. 11. 1893}\newcommand{\editorInnen}{Herausgegeben von Jahnke, SelmaMüller, Martin Anton}\input{../tex-inputs/latex-leseansicht-abspann}
